\documentclass[11pt]{article}
\usepackage[margin=1in]{geometry}
\usepackage[T1]{fontenc}
\usepackage{lmodern,amsmath,amssymb,amsthm,microtype}
\usepackage[colorlinks=true,urlcolor=blue!50!black,linkcolor=black]{hyperref}
\usepackage{xcolor}
\newtheorem{theorem}{Theorem}
\newtheorem{lemma}[theorem]{Lemma}
\newcommand{\fl}[1]{\lfloor #1\rfloor}
\newcommand{\ce}[1]{\lceil #1\rceil}
\hypersetup{pdftitle={Minimizing a mixed star in a tournament}}
\title{\bfseries Minimizing a mixed star in a tournament}
\author{}
\date{Working note \(\cdot\) November 2026}
\begin{document}
\maketitle
\begin{abstract}\noindent
How few stars with two arrows entering the center and one leaving it can
an \(n\)-vertex tournament contain? An optimal construction puts a nearly
regular tournament above a transitive tournament on
\(\fl{(n-1)/2}\) vertices. Its count is
\((31/768)n^4+O(n^3)\) when the two incoming leaves are unordered.
The same degree comparison gives an exact finite optimization rule for
every cubic function of tournament scores.
\end{abstract}

\section*{The construction and the count}
Put \(N=n-1\). For \(0\le b\le N\), take a transitive tournament on \(b\)
vertices, and direct every edge from the remaining \(n-b\) vertices
towards these \(b\). Make the upper tournament regular if its order is
odd, and nearly regular (internal degrees differing by one) otherwise.
Odd regular tournaments exist by the cyclic construction; deleting one
vertex of the next odd regular tournament gives the even case.
Write \(D_b\) for the resulting multiset of outdegrees. It consists of
\(0,\ldots,b-1\), followed by \(n-b\) values with mean \(v=(b+N)/2\).
They all equal \(v\) when it is integral, and otherwise half equal
\(\fl v\) and half \(\ce v\).

\begin{theorem}\label{th:star}
For every \(n\ge4\), the minimum number of mixed stars with one outgoing
leaf and two unordered incoming leaves is
\[
   \sum_{x\in D_b}x\binom{N-x}{2},\qquad b=\fl{N/2}.
\]
Edges among the leaves are unrestricted. Reversing every edge gives the
same minimum for the opposite orientation.
\end{theorem}

A vertex of outdegree \(x\) contributes exactly \(x\binom{N-x}{2}\)
choices. For \(n=7\) the construction has degrees \(0,1,2,4,4,5,5\)
and gives 30 stars. Minimizing the number of possible centers gives a
different answer: the three-center hosts have at least 31 stars, as the
separate exhaustive seven-vertex calculation confirms.

On dividing by \(n^4\), the lower degrees give a Riemann sum and the
upper half concentrate at normalized degree \(3/4\). The limiting count is
\[
 \frac12\left(\int_0^{1/2}x(1-x)^2\,dx
       +\tfrac12\cdot\tfrac34(1-\tfrac34)^2\right)=\frac{31}{768}.
\]
The rounding of the cutoff and of the upper degrees contributes \(O(n^3)\)
to the unnormalized count. Labeling the two incoming leaves doubles it.

\newpage
\section*{A lower bound with equality}
For a multiset \(D\) of tournament degrees write
\(S_f(D)=\sum_{x\in D} f(x)\). The next observation is the elementary
majorization inequality for scores, included to specify its direction.
\begin{lemma}\label{lem:major}
For concave \(g\) on \([0,N]\),
\(S_g(D)\ge\sum_{i=0}^{N} g(i)\).
\end{lemma}
\begin{proof}
For the \(j\) smallest scores, the sum is at least \(\binom j2\) by
counting edges among those vertices. The full sum is \(\binom n2\).
Thus for each integer \(t\), \(\sum_{x\in D}(t-x)_+\) is at most
\(\max_{0\le j\le n}(jt-\binom j2)=\sum_{i=0}^{N}(t-i)_+\).
On the integer grid a concave function is an affine function minus a
nonnegative combination of these hinges: the coefficient at \(t\) is
\(2g(t)-g(t-1)-g(t+1)\), for \(1\le t<N\). Summing proves the claim.
\end{proof}

\begin{lemma}\label{lem:cert}
Let \(f_s(x)=x^3-sx^2\), with integer \(N\le s\le3N\). Put
\[
 b=\ce{(s-N)/2},\quad u=s-N-b,\quad
 p=\fl{(N+b)/2},\quad q=\ce{(N+b)/2}.
\]
Then \(D_b\) minimizes \(S_{f_s}\).
\end{lemma}
\begin{proof}
We have \(u=b-1\) or \(u=b\), and \(0\le u\le p\le q\).
Since \(p+q=N+b\), the polynomial
\[
   L(x)=f_s(x)-(x-u)(x-p)(x-q)
\]
is affine. Define \(g=f_s\) up to \(u\) and \(g=L\) from \(u\) onward.
On the left \(f_s''(x)=6x-2s\le0\) since \(s=u+p+q\ge3u\).
At the join \(f_s'(u)-L'(u)=(u-p)(u-q)\ge0\), so \(g\) is concave.
For integer \(x\ge u\), the factors \(x-p\) and \(x-q\) have
nonnegative product because \(q-p\le1\). Thus \(g(x)\le f_s(x)\)
at every integer score, including the rounding case \(p\ne q\).
Lemma~\ref{lem:major} gives
\[
 S_{f_s}(D)\ \ge\ S_g(D)\ \ge\ \sum_{i=0}^N g(i).
\]
For \(D_b\) both inequalities are equalities. Every score either lies
at or below \(u\), or equals \(p\) or \(q\). Moreover replacing
\(b,\ldots,N\) by the upper degrees preserves both their number and sum,
and \(g\) is affine on that interval. The lower degrees are unchanged.
\end{proof}

\noindent\emph{Proof of Theorem~\ref{th:star}.}
Twice the cost at \(x\) is \(f_{2N-1}(x)+N(N-1)x\). The last term has
constant total. Lemma~\ref{lem:cert} gives
\(b=\ce{(N-1)/2}=\fl{N/2}\).
\hfill\(\square\)

\newpage
\section*{Arbitrary cubic objectives}
\begin{theorem}
For any real \(s\) and \(n\ge2\), the minimum of \(S_{f_s}\) is the
least of its values on \(D_0,\ldots,D_N\). Consequently any cubic
polynomial of the scores with positive leading coefficient attains its
minimum in this list. With negative leading coefficient the same holds
after reversing all arcs.
\end{theorem}
\begin{proof}
For \(s\le N\), take \(b=0,\ u=s-N\le0,\ p=\fl{N/2},\
q=\ce{N/2}\). The same affine \(L\) lies below \(f_s\) at each integer
in \([0,N]\), with equality at \(p,q\). The affine totals agree for all
tournaments, so \(D_0\) attains the minimum. For \(s\ge3N\), \(f_s\) is
concave throughout \([0,N]\), and Lemma~\ref{lem:major} gives \(D_N\).

Between consecutive integers in \([N,3N]\), it suffices to find one
construction optimal at both endpoints: every difference between two
constructions' objective values is affine in \(s\), as is the difference
with any other tournament. Write the left endpoint as \(N+j\). If \(j\)
is odd, Lemma~\ref{lem:cert} chooses the same \(b\) at both endpoints.
If \(j=2b\), the two prescribed prefixes have lengths \(b,b+1\).
At both endpoints the break \(u\) in the certificate equals \(b\).
Write \(v=(N+b)/2\).
\begin{enumerate}
\item If \(v\) is an integer, use \(D_b\). Its upper score \(v\) is the
lower contact integer for \(D_{b+1}\), whose upper mean is \(v+1/2\).
\item If \(v\) is a half-integer, use \(D_{b+1}\). Its extra lower
score \(b\) is the break itself and its upper score \(v+1/2\) is the
upper contact integer for \(D_b\).
\end{enumerate}
In either case the alternative construction also gives equality in both
inequalities of the certificate, because the replacement above \(b\)
preserves the sum on the affine part. This proves interpolation.
Constant and linear terms have fixed totals. For a negative leading
coefficient replace every degree \(x\) by \(N-x\).
\end{proof}

\paragraph{Checks and context.}
Independent score enumeration checks the cubic minimum on a grid of
slopes spaced \(1/3\) through \(n=10\); a score-sequence optimization
checks the star count through \(n=24\).
The majorization of tournament scores is classical. Th\"ornblad studies
the corresponding limiting degree distributions
(\href{https://arxiv.org/abs/1611.09579}{arXiv:1611.09579}).
Fox, Himwich, Mani and Zhou classify which star orientations have
asymptotic densities maximized or minimized by random tournaments
(\href{https://arxiv.org/abs/2402.08418}{Theorem 1.11}).
The present calculation concerns the opposite extremum for this mixed
star. Pooling an upper interval using an affine bound also occurs in
the continuous optimization of Kolotilin, Mylovanov and Zapechelnyuk
(\href{https://doi.org/10.3982/TE4071}{Theoretical Economics, 2022}).
The scope of prior finite integer-score results remains to be established.
\end{document}
